\section*{Chapter \ref{ch:b1:organomagnesium} --- Organomagnesium Reagents}

\begin{solution}{exo:b1:organomagnesium:1}
\ce{CH3CH2Br + Mg -> CH3CH2MgBr}; \ce{C6H5Br + Mg -> C6H5MgBr} (in dry
ether). C--Br: Br (2.96) more electronegative than C (2.55), carbon
positive. C--Mg: Mg (1.31) less electronegative, carbon negative.
\end{solution}

\begin{solution}{exo:b1:organomagnesium:2}
Water: \ce{CH3MgI + H2O -> CH4 + Mg(OH)I}.
Ethanol: \ce{CH3MgI + CH3CH2OH -> CH4 + CH3CH2OMgI}.
Ethanoic acid: \ce{CH3MgI + CH3COOH -> CH4 + CH3COOMgI}.
Deuterium oxide: \ce{CH3MgI + D2O -> CH3D + Mg(OD)I}.
\end{solution}

\begin{solution}{exo:b1:organomagnesium:3}
Methanal: propan-1-ol (primary). Ethanal: butan-2-ol (secondary).
Propanone: 2-methylbutan-2-ol (tertiary). Benzaldehyde:
1-phenylpropan-1-ol (secondary). Carbon dioxide: propanoic acid.
\end{solution}

\begin{solution}{exo:b1:organomagnesium:4}
Its OH group is acidic enough to destroy the reagent as soon as it forms:
the molecule would protonate itself (or its neighbours), giving the
magnesium alkoxide of butan-1-ol, \ce{CH3(CH2)3OMgBr}, not a usable
reagent. The OH must first be protected
(\cref{ch:b1:acetals-protection} shows the idea for an aldehyde).
\end{solution}

\begin{solution}{exo:b1:organomagnesium:5}
The C--Mg pair attacks the carbonyl carbon of ethanal, the C=O $\pi$ pair
going to O: \ce{CH3CH(CH3)O^-} with \ce{MgBr+}; then \ce{H3O+} gives a
proton to the oxygen: propan-2-ol. Propan-2-ol carries two methyl groups
on the carbinol carbon: it is \omterm{def:b1:stereochemistry-in-depth:stereocentre}{achiral}. (With a different R, e.g.
ethylmagnesium bromide on ethanal, the product butan-2-ol is \omterm{def:b1:stereochemistry-in-depth:stereocentre}{chiral} but
racemic: the planar carbonyl is attacked on both faces equally.)
\end{solution}

\begin{solution}{exo:b1:organomagnesium:6}
The carbinol carbon carries one methyl and two ethyl groups:
ethylmagnesium bromide with butan-2-one, or methylmagnesium bromide with
pentan-3-one.
\end{solution}

\begin{solution}{exo:b1:organomagnesium:7}
Prepare \ce{(CH3)2CHMgBr} from 2-bromopropane and magnesium in dry ether,
pour it onto solid \ce{CO2}, then acidify: \ce{(CH3)2CHCOOH}. $n = 12.3/123.0 = 0.100$ mol;
$0.100 \times 0.70 \times 88.0 = \qty{6.2}{g}$.
\end{solution}

\begin{solution}{exo:b1:organomagnesium:8}
$0.18/18.0 = \qty{0.010}{mol}$ of water destroys \qty{0.010}{mol} of
reagent: \qty{20}{\percent}; benzene forms, \ce{C6H5MgBr + H2O -> C6H6 +
Mg(OH)Br}.
\end{solution}

\begin{solution}{exo:b1:organomagnesium:9}
The reagent, a \omterm{def:b1:electronic-effects:nucleophile}{nucleophile}, reacts with bromobenzene still present:
overall \ce{C6H5MgBr + C6H5Br -> C6H5-C6H5 + MgBr2}. Adding the bromide
slowly to excess magnesium keeps its concentration low: it meets fresh
metal rather than the reagent already formed.
\end{solution}

\begin{solution}{exo:b1:organomagnesium:10}
1-Phenylethan-1-ol: \ce{CH3MgBr} (from bromomethane) with benzaldehyde,
or \ce{C6H5MgBr} (from bromobenzene) with ethanal. 2-Phenylpropan-2-ol:
\ce{C6H5MgBr} with propanone, or \ce{CH3MgBr} with
1-phenylethan-1-one.
\end{solution}

\begin{solution}{exo:b1:organomagnesium:11}
The ketone is added to the reagent so that the reagent, prepared in the
dry flask, never leaves it and the heat released by each portion is
absorbed by the boiling ether. Water alone would give the gelatinous
basic salt \ce{Mg(OH)Br}, which traps the product; cold dilute acid
dissolves it as \ce{Mg^2+}, and the cold limits side reactions of the
tertiary alcohol in acid.
\end{solution}

\begin{solution}{exo:b1:organomagnesium:12}
Each mole of reagent gives with water one mole of basic magnesium salt,
\ce{RMgBr + H2O -> RH + Mg(OH)Br}, whose \ce{OH-} is titrated by the acid:
the reagent amount equals the acid used, \qty{0.088}{mol}, a \omterm{def:b1:extent-q-and-k:fractional-extent}{yield} of
\qty{88}{\percent}.
\end{solution}

\begin{solution}{pb:b1:organomagnesium:1}
\textbf{1.} \ce{C6H5Br + Mg -> C6H5MgBr}.
\textbf{2.} Mg $0.535/24.3 = \qty{0.0220}{mol}$; \ce{C6H5Br} $3.14/156.9 =
\qty{0.0200}{mol}$: magnesium in slight excess.
\textbf{3.} Water destroys the reagent: \ce{C6H5MgBr + H2O -> C6H6 +
Mg(OH)Br}.
\textbf{4.} It stops the water vapour of the air entering through the
condenser.
\textbf{5.} The ether gives \omterm{def:b1:lewis-resonance-vsepr:lewis}{lone pairs} to magnesium (\omterm{def:b1:lewis-resonance-vsepr:lewis-acid}{Lewis base}):
without it the reagent does not form or stay in solution.
\textbf{6.} \ce{C6H5MgBr + C6H5Br -> C6H5-C6H5 + MgBr2}; biphenyl
$0.15/154.0 = \qty{9.7e-4}{mol}$, which consumed \qty{9.7e-4}{mol} of
bromobenzene and as much reagent: \qty{1.9e-3}{mol} of bromobenzene lost
in all.
\textbf{7.} $3.28/182.0 = \qty{0.0180}{mol}$.
\textbf{8.} The C--Mg pair attacks the carbonyl carbon, the $\pi$ pair
goes to oxygen: \ce{(C6H5)3C-O^-} \ce{MgBr+}.
\textbf{9.} Reagent available at most $0.0200 - 0.0019 = \qty{0.0181}{mol}$;
benzophenone \qty{0.0180}{mol}: benzophenone limits, just.
\textbf{10.} No: the carbinol carbon carries three identical phenyl groups.
\textbf{11.} The reagent stays in its dry flask, and the heat of each
portion of ketone is taken up by the boiling ether.
\textbf{12.} It would destroy part of the reagent (benzene) and leave
benzophenone unreacted.
\textbf{13.} \ce{(C6H5)3COMgBr + H3O+ -> (C6H5)3COH + Mg^2+ + Br- + H2O}.
\textbf{14.} The acid dissolves the basic magnesium salts that water
alone would precipitate; the cold limits the formation of the
\omterm{def:b1:electronic-effects:intermediates}{carbocation} of question 18.
\textbf{15.} Triphenylmethanol and biphenyl in the ether layer; the
magnesium salts in the aqueous layer.
\textbf{16.} Washing (or triturating) the solid with cold light petroleum:
biphenyl dissolves, triphenylmethanol stays.
\textbf{17.} A sharp melting point equal to the tabulated one (or a single
spot in thin-layer chromatography).
\textbf{18.} \ce{(C6H5)3COH + H+ -> (C6H5)3C+ + H2O}: a tertiary \omterm{def:b1:electronic-effects:intermediates}{carbocation}
whose charge is spread over three benzene rings by resonance, very
stable, and coloured because of that extended conjugation.
\textbf{19.} $3.75/260.0 = \qty{0.0144}{mol}$.
\textbf{20.} $0.0144/0.0180 = \qty{80}{\percent}$.
\textbf{21.} Losses in transfers and recrystallisation; reagent destroyed
by traces of moisture; incomplete reaction.
\textbf{22.} Benzoic acid, \ce{C6H5COOH}, after acidification.
\textbf{23.} $0.0181 \times 122.0 = \qty{2.21}{g}$.
\textbf{24.} \ce{C6H5MgBr + CO2 -> C6H5COOMgBr}, then
\ce{C6H5COOMgBr + H3O+ -> C6H5COOH + Mg^2+ + Br- + H2O}.
\textbf{25.} \textbf{\qty{80}{\percent}}.
\end{solution}
